% =============================================================================
% 3.2 Architecture technique
% =============================================================================
\section{Architecture technique}
\label{sec:ch3-technique}

Ce chapitre nomme les composants qui incarnent les rôles définis en
méthodologie. Le tableau~\ref{tab:composants} en donne la correspondance.

\begin{btptable}{Correspondance rôles méthodologiques --- composants techniques}
\label{tab:composants}
\begin{tabularx}{\linewidth}{@{}L{4.0cm} Y L{3.6cm}@{}}
\toprule
\rowcolor{btppale}
\btpth{Rôle} & \btpth{Composant} & \btpth{Justification brève} \\
\midrule
Interface applicative & FastAPI & Contrats HTTP explicites, validation Pydantic. \\
File de travaux & ARQ et Redis & Ingestion asynchrone, reprises. \\
Analyse de documents & Docling / PaddleOCR & Numérique vs numérisé. \\
Indexation et recherche & LlamaIndex et Qdrant & Hybride dense et creux natif. \\
Plongements & BGE-M3 & Mixte, multilingue, français inclus. \\
Ré-ordonnancement & BGE-reranker-v2-m3 & Modèle croisé, couple (requête, extrait). \\
Orchestration & LangGraph & Graphes d'étapes et interruptions HITL. \\
Modèle de langage & Ollama (local) ou API & Exécution locale prioritaire. \\
Faits métier & PostgreSQL & Versions, citations, revues. \\
Fichiers & MinIO & Originaux et exports. \\
Graphe (phase 5) & Neo4j & Reporté. \\
\bottomrule
\end{tabularx}
\end{btptable}

\subsection{Choix d'extraction documentaire}

Deux outils coexistent, conformément à la bifurcation identifiée à la
section~\ref{sec:ch1-documents}. Docling est retenu pour les PDF
numériques, car il préserve titres, listes et tableaux. PaddleOCR, avec
analyse de mise en page, est retenu pour les scans. Le basculement est
automatique, fondé sur la détection d'une couche texte extractible.

\subsection{Choix d'indexation}

Qdrant est retenu pour trois propriétés alignées sur le cahier des charges
de la section~\ref{sec:ch1-vecteurs}~: filtrage par métadonnées avant
similarité, représentation mixte dense et creuse, identifiants stables
rattachables aux extraits stockés dans PostgreSQL. BGE-M3 produit
conjointement les deux représentations, ce qui évite d'entretenir deux
modèles de plongement.

\subsection{Choix d'orchestration}

LangGraph est retenu parce qu'il expose nativement des graphes d'étapes
interruptibles. Le point d'interruption de revue n'est pas un rappel
applicatif greffé après coup~: il est un n\oe ud du graphe, ce qui
garantit qu'aucune section n'est persistée comme approuvée sans décision
humaine. LlamaIndex fournit les abstractions de récupération (enfants,
parents, fusion, ré-ordonnancement) sans dicter l'orchestration.

\begin{btpnote}[Substitutabilité]
Ces choix sont des incarnations, pas des fondements. Remplacer Qdrant par
un autre magasin hybride, ou Ollama par une API, ne change pas le
workflow. C'est précisément le bénéfice de l'architecture en couches.
\end{btpnote}
